% Four physical strands. Read operators from bottom (input) to top (output).
% An ordinary swap is TWO separate elbows, never a four-way junction.
% Circular boundaries denote local L(0) tensors, including their scalar.
% A superscript inverse on the boundary denotes the inverse local tensor.
\newcommand{\xxzTensorBoundary}[2]{\draw (#1,#2) circle (.31);}
\newcommand{\xxzSwap}[3]{%
 \xxzTensorBoundary{#1}{#2}
 \ifnum#3=0 \node[font=\Large,anchor=south west,inner sep=0pt] at (#1+.21,#2+.23) {$-1$};\fi
 \ifnum#3=1
  \draw (#1,#2-.43)--(#1,#2-.12)--(#1-.12,#2)--(#1-.43,#2);
  \draw (#1+.43,#2)--(#1+.12,#2)--(#1,#2+.12)--(#1,#2+.43);
 \else
  \draw (#1,#2-.43)--(#1,#2-.12)--(#1+.12,#2)--(#1+.43,#2);
  \draw (#1-.43,#2)--(#1-.12,#2)--(#1,#2+.12)--(#1,#2+.43);
 \fi}
\newcommand{\xxzDerivative}[2]{%
 \draw (#1,#2-.43)--(#1,#2+.43);
 \draw (#1-.43,#2)--(#1+.43,#2);
 \draw[fill=white] (#1,#2) circle (.34);
 \node at (#1,#2) {$L'$};}
\newcommand{\xxzRow}[3]{%
 % row y, orientation, derivative site (0 means no insertion)
 \foreach \j in {1,...,4}{%
  \pgfmathsetmacro{\xx}{1.25*(\j-1)}
  \ifnum\j=#3 \xxzDerivative{\xx}{#1}\else\xxzSwap{\xx}{#1}{#2}\fi}
 \foreach \j in {0,...,2}{\draw (1.25*\j+.43,#1)--(1.25*\j+.82,#1);}}
% Matching dashed cuts indicate the same periodic auxiliary space.
\newcommand{\xxzPeriodicEnds}[1]{%
 \draw (-.43,#1)--(-.72,#1);
 \draw (4.18,#1)--(4.47,#1);
 \draw[densely dashed,line width=.6pt] (-.72,#1-.62)--(-.72,#1+.62);
 \draw[densely dashed,line width=.6pt] (4.47,#1-.62)--(4.47,#1+.62);}
\newcommand{\xxzSiteLabels}[1]{%
 \foreach \j in {1,...,4}{\node at ({1.25*(\j-1)},#1) {$\j$};}}
